\section{Random features}
\label{sec:rf}

Keep the predictor, replace the infinite map by \(2Q\) numbers with \(2Q<n\).
Draw frequencies \(\omega_q\sim\mathcal{N}(0,I_d)\) once, before seeing data, and set
\begin{equation}
\psi(x)
= \frac1{\sqrt Q}
\begin{pmatrix}
\cos(\omega_1^\top x) \\
\sin(\omega_1^\top x) \\
\vdots \\
\cos(\omega_Q^\top x) \\
\sin(\omega_Q^\top x)
\end{pmatrix}
\in\R^{2Q}.
\label{eq:psi}
\end{equation}
Rows \(\psi(x_i)^\top\) form \(\Psi\in\R^{n\times 2Q}\). The estimator is ordinary ridge regression in these coordinates,
\begin{equation}
\hat f(x_*)
= \psi(x_*)^\top(\Psi^\top\Psi+\lambda I)^{-1}\Psi^\top y.
\label{eq:rf}
\end{equation}
Cost \(O(nQd+nQ^2+Q^3)\), linear in the sample size.

It matches kernel ridge when \(\psi(x)^\top\psi(x')\) matches \(k(x,x')\).
Bochner's theorem writes the Gaussian kernel as a Fourier expectation \cite{bochner1933,rudin1962}.
Rahimi and Recht turn that expectation into features \cite{rahimi2007}.
Because \(k\) is real and even, the expectation is an average of cosines,
\begin{equation}
k(x,x')
= \kappa(x-x')
= \E_{\omega\sim\mathcal{N}(0,I_d)}
\bigl[\cos\bigl(\omega^\top(x-x')\bigr)\bigr].
\label{eq:bochner}
\end{equation}
The average of \(Q\) independent draws is
\begin{equation}
k(x,x')
\approx \frac1Q\sum_{q\le Q}\cos\bigl(\omega_q^\top(x-x')\bigr)
= \psi(x)^\top\psi(x').
\label{eq:lln}
\end{equation}
The cosine splits by the angle formula, which is why \(\psi\) stores both \(\cos(\omega_q^\top x)\) and \(\sin(\omega_q^\top x)\).
The weak law puts the error near \(1/\sqrt{2Q}\). At \(Q=5000\) that is about \(0.01\).

Then \(K\approx\Psi\Psi^\top\) and \(k_*\approx\Psi\psi(x_*)\).
The push-through identity
\begin{equation}
\Psi^\top(\Psi\Psi^\top+\lambda I)^{-1}
= (\Psi^\top\Psi+\lambda I)^{-1}\Psi^\top
\label{eq:push}
\end{equation}
turns the \(n\times n\) kernel predictor into the \(2Q\times 2Q\) estimator \eqref{eq:rf}.

\begin{table}[ht]
\centering
\caption{Three estimators of the same squared loss.}
\label{tab:costs}
\begin{tabular}{lll}
\toprule
Estimator & Coordinates & Work \\
\midrule
Linear ridge & \(x\in\R^d\) & \(O(nd^2+d^3)\) \\
Kernel ridge & \(\varphi(x)\in\R^\infty\) & \(O(n^2 d+n^3)\) \\
Random features & \(\psi(x)\in\R^{2Q}\) & \(O(nQd+nQ^2+Q^3)\) \\
\bottomrule
\end{tabular}
\end{table}

\noindent
The kernel buys every scaled monomial. Random features buy the same predictor back at a cost linear in \(n\).
